% !TeX root = ../presentacion.tex
% ===========================================================================
%  7. Reglas de negocio — 1:10 (dos láminas)
%
%  Criterio: «Reglas de negocio y validaciones» (10 %), donde la rúbrica
%  señala «en especial la validación de solapamiento de horarios». Por eso
%  RN-02 se lleva una lámina entera.
%
%  Las clases y métodos de la tabla están verificados contra el código:
%  BookingService.crear, Booking.sePuedeCancelarEn, Booking.cancelar,
%  BookingService.verificarTope, Booking.estadoVigente, CourtService.
%  El fragmento SQL es literal de
%  backend/ms-reservas/src/main/resources/db/migration/V1__reservas.sql
% ===========================================================================

\bloque{Reglas de negocio}{1:10}{\exponeBackend}%
       {Reglas de negocio y validaciones (10\,\%)}

\begin{frame}{Ocho reglas, y dónde vive cada una}
  \note{No leer la tabla entera. Señalar la columna de la derecha y decir la
  idea: ninguna regla vive solo en la prosa del informe, y ninguna fila está
  sin prueba.}

  \begin{center}
    \scriptsize
    \setlength{\tabcolsep}{4pt}
    \renewcommand{\arraystretch}{1.08}
    \begin{tabular}{@{}lll@{}}
      \toprule
      & \textbf{Regla} & \textbf{Dónde vive}\\
      \midrule
      \rn{01} & Cancha, fecha y bloque de una hora         & \id{BookingService.crear} · \id{TimeSlot}\\
      \rn{02} & Sin solapamiento, ni bajo concurrencia     & Servicio \alert{y} \id{EXCLUDE} en la base\\
      \rn{03} & Cada quien cancela lo suyo; el admin, todo & \id{BookingService.cancelar}\\
      \rn{04} & No se cancela una reserva ya empezada      & \id{Booking.sePuedeCancelarEn}\\
      \rn{05} & Cancelar libera el bloque, sin borrar      & \id{Booking.cancelar}\\
      \rn{06} & Tope de reservas activas (3 por defecto)   & \id{BookingService.verificarTope}\\
      \rn{07} & Solo el admin gestiona el catálogo         & \id{CourtService} en \id{ms-canchas}\\
      \rn{08} & Confirmada / Cancelada / Finalizada        & \id{Booking.estadoVigente}\\
      \bottomrule
    \end{tabular}
  \end{center}

  \vspace{2pt}
  {\scriptsize Las ocho están cubiertas por \textbf{73 pruebas
  automatizadas} en 12 clases. La correspondencia completa, en el informe
  (\ent{1}, §7).}
\end{frame}

\begin{frame}[fragile]{\rn{02}: validar en el servicio no basta}
  \note{La lámina técnica de la charla. Contar la carrera: dos peticiones
  simultáneas preguntan «¿está libre?», las dos oyen que sí, las dos
  insertan. La comprobación en el servicio es necesaria para dar un mensaje
  decente; la restricción es la que hace que la regla sea CIERTA.}

  \begin{columns}[T]
    \begin{column}{0.44\textwidth}
      \begin{block}{La carrera}
        \scriptsize
        Dos peticiones sobre el mismo bloque consultan antes de que ninguna
        haya escrito. Ambas lo ven libre. Ambas insertan.
      \end{block}
      \begin{block}{Quién impone la regla}
        \scriptsize
        PostgreSQL. Si la carrera ocurre, la segunda inserción falla y el
        servicio la traduce a \id{409 Conflict}.
      \end{block}
      {\scriptsize\color{textoSuave} \id{ConcurrenciaReservaIT} lanza las dos
      a la vez, 10 veces seguidas, y exige que gane exactamente una.}
    \end{column}
    \begin{column}{0.53\textwidth}
      \begin{lstlisting}[language=sql-pg,basicstyle=\ttfamily\tiny]
ALTER TABLE reserva
  ADD CONSTRAINT reserva_sin_solape
  EXCLUDE USING gist (
    cancha_id WITH =,
    tsrange(fecha + hora_inicio,
            fecha + hora_fin, '[)') WITH &&
  ) WHERE (estado = 'CONFIRMADA');
      \end{lstlisting}
      {\scriptsize\color{textoSuave}
      \id{WHERE estado = 'CONFIRMADA'} es lo que hace que cancelar libere el
      bloque (\rn{05}) sin borrar la fila, conservando la traza de \rn{08} y
      el dato del reporte de cancelaciones.}
    \end{column}
  \end{columns}
\end{frame}
